% Generated. Do not edit.
% Projection scriptural-date-location/section, version 1: the Scriptural Date and
% Location page, generated whole from research/chronology.toml, the Scripture
% chronology corpus and the source library by tools/tpt document-types refresh.
% document: liturgy/roman-rite/1962/propers/temporal/60-twentieth-after-pentecost
\runninghead{Scriptural Date and Location}
\section{Scriptural Date and Location}
\label{sec:date-location}

{\small Each passage the Mass reads or sings from Scripture, in canonical order. The Context column names what is dated --- the event the passage narrates, a setting it remembers, the traditional attribution of its author, or the composition of its book --- and where it belongs; the Date column gives the answer of the Scripture chronology corpus for that subject, with every alternative it records. Texts of the Mass that are not drawn from Scripture have no row. The evidence for each date follows the table.\par}
\medskip
% Generated. Do not edit.
% Deterministic projection of research/chronology.toml.
% schema: 5
% document: liturgy/roman-rite/1962/propers/temporal/60-twentieth-after-pentecost
% generated-by: tools/tpt proper-chronology annotations
\newcommand{\chronologyannotationclaim}[6]{#6}
\newcommand{\chronologyannotationcomparisonclaim}[8]{#8}
\newcommand{\chronologyannotationanchorclaim}[9]{#9}
\newcommand{\chronologyannotationanchorgroup}[8]{#8}
\newcommand{\chronologyannotationreach}[3]{}
\newcommand{\chronologyannotationgroup}[3]{#3}
\newcommand{\chronologyannotationcomparisongroup}[4]{#4}
\newcommand{\chronologyannotation}[1]{%
  \ifcsname triptychchronologyannotation@#1\endcsname
    \csname triptychchronologyannotation@#1\endcsname
  \else
    \PackageError{triptych}{No chronology annotation for #1}{}%
  \fi}
% comparison-dependency: src/claude/liturgy/roman-rite/1962/propers/temporal/60-twentieth-after-pentecost/research/chronology-profile-comparisons.toml sha256=139a40c7f50531ff34337f64368f66dc23e9909ef9d06f2802237a8200eee01c
% introit: Introit
\expandafter\def\csname triptychchronologyannotation@introit\endcsname{%
  \chronologyannotationgroup{historical-setting}{nonuniform}{\textbf{Historical setting} -- No single occasion date applies across every cited locus.}%
  \space \chronologyannotationgroup{composition}{nonuniform}{\textbf{Composition} -- No single date applies across every cited locus.}%
}
% epistle: Epistle
\expandafter\def\csname triptychchronologyannotation@epistle\endcsname{%
  \chronologyannotationgroup{composition}{disputed}{\textbf{Composition} -- disputed: \chronologyannotationreach{Eph.5.15}{Eph}{true}\chronologyannotationreach{Eph.5.16}{Eph}{true}\chronologyannotationreach{Eph.5.17}{Eph}{true}\chronologyannotationreach{Eph.5.18}{Eph}{true}\chronologyannotationreach{Eph.5.19}{Eph}{true}\chronologyannotationreach{Eph.5.20}{Eph}{true}\chronologyannotationreach{Eph.5.21}{Eph}{true}\chronologyannotationclaim{composition.epistle-to-the-ephesians}{composition}{catholic-traditional-v1}{disputed}{a period between 58 and 63}{A.D. 58--63}; \chronologyannotationreach{Eph.5.15}{Eph}{true}\chronologyannotationreach{Eph.5.16}{Eph}{true}\chronologyannotationreach{Eph.5.17}{Eph}{true}\chronologyannotationreach{Eph.5.18}{Eph}{true}\chronologyannotationreach{Eph.5.19}{Eph}{true}\chronologyannotationreach{Eph.5.20}{Eph}{true}\chronologyannotationreach{Eph.5.21}{Eph}{true}\chronologyannotationclaim{composition.epistle-to-the-ephesians}{composition}{catholic-traditional-v1}{disputed}{(Philemon; Colossians; Ephesians; Philippians), 61}{A.D. 61}.}%
  \space \chronologyannotationcomparisongroup{catholic-critical-v1}{composition}{composition-only}{\textbf{Composition (Catholic critical chronology, for comparison)} -- disputed: \chronologyannotationreach{Eph.5.15}{Eph}{true}\chronologyannotationreach{Eph.5.16}{Eph}{true}\chronologyannotationreach{Eph.5.17}{Eph}{true}\chronologyannotationreach{Eph.5.18}{Eph}{true}\chronologyannotationreach{Eph.5.19}{Eph}{true}\chronologyannotationreach{Eph.5.20}{Eph}{true}\chronologyannotationreach{Eph.5.21}{Eph}{true}\chronologyannotationcomparisonclaim{critical.ephesians-later-disciple}{composition}{catholic-critical-v1}{catholic-critical-v1}{disputed}{passage.united-states-conference-of-catholic-bishops.new-american-bible-revised-edition.english-usccb-web-2026-09-28.ephesians-introduction}{around A.D. 80--100}{Ephesians under the later-disciple hypothesis (approximate), c. A.D. 80--100}.}%
}
% gradual: Gradual
\expandafter\def\csname triptychchronologyannotation@gradual\endcsname{%
  \chronologyannotationgroup{traditional-attribution}{disputed}{\textbf{Traditional attribution} -- disputed: \chronologyannotationreach{Ps.144.15}{Ps.107 Ps.144}{true}\chronologyannotationreach{Ps.144.16}{Ps.107 Ps.144}{true}\chronologyannotationclaim{israel.monarchy.david-traditional-era}{traditional-attribution}{catholic-traditional-v1}{disputed}{reigned from 1055 to 1015 B.C.}{David (reign in the usual chronology), B.C. 1055--1015}.}%
  \space \chronologyannotationgroup{historical-setting}{research-pending}{\textbf{Historical setting} -- Occasion date unresolved.}%
  \space \chronologyannotationgroup{composition}{preferred}{\textbf{Composition}: \chronologyannotationreach{Ps.144.15}{Ps}{true}\chronologyannotationreach{Ps.144.16}{Ps}{true}\chronologyannotationclaim{critical.psalms.latest-composition-boundary}{composition}{catholic-critical-v1}{preferred}{before the Maccabean period, around 165 B.C.; no individual psalm can be dated securely}{Before c. 165 B.C.}}%
}
% alleluia: Alleluia
\expandafter\def\csname triptychchronologyannotation@alleluia\endcsname{%
  \chronologyannotationgroup{traditional-attribution}{disputed}{\textbf{Traditional attribution} -- disputed: \chronologyannotationreach{Ps.107.2}{Ps.107 Ps.144}{true}\chronologyannotationclaim{israel.monarchy.david-traditional-era}{traditional-attribution}{catholic-traditional-v1}{disputed}{reigned from 1055 to 1015 B.C.}{David (reign in the usual chronology), B.C. 1055--1015}.}%
  \space \chronologyannotationgroup{historical-setting}{research-pending}{\textbf{Historical setting} -- Occasion date unresolved.}%
  \space \chronologyannotationgroup{composition}{preferred}{\textbf{Composition}: \chronologyannotationreach{Ps.107.2}{Ps}{true}\chronologyannotationclaim{critical.psalms.latest-composition-boundary}{composition}{catholic-critical-v1}{preferred}{before the Maccabean period, around 165 B.C.; no individual psalm can be dated securely}{Before c. 165 B.C.}}%
}
% gospel: Gospel
\expandafter\def\csname triptychchronologyannotation@gospel\endcsname{%
  \chronologyannotationgroup{narrated-event}{preferred}{\textbf{Event}: \chronologyannotationreach{John.4.46}{John.4.46-53}{false}\chronologyannotationreach{John.4.47}{John.4.46-53}{false}\chronologyannotationreach{John.4.48}{John.4.46-53}{false}\chronologyannotationreach{John.4.49}{John.4.46-53}{false}\chronologyannotationreach{John.4.50}{John.4.46-53}{false}\chronologyannotationreach{John.4.51}{John.4.46-53}{false}\chronologyannotationreach{John.4.52}{John.4.46-53}{false}\chronologyannotationreach{John.4.53}{John.4.46-53}{false}\chronologyannotationclaim{life-of-christ.healing-of-the-rulers-son-at-cana}{narrated-event}{catholic-traditional-v1}{preferred}{Passover, A.U.C. 779 - about Pentecost, 780}{The healing of the ruler's son at Cana, A.D. 26--27 (derived)}.}%
  \space \chronologyannotationgroup{composition}{disputed}{\textbf{Composition} -- disputed: \chronologyannotationreach{John.4.46}{John}{true}\chronologyannotationreach{John.4.47}{John}{true}\chronologyannotationreach{John.4.48}{John}{true}\chronologyannotationreach{John.4.49}{John}{true}\chronologyannotationreach{John.4.50}{John}{true}\chronologyannotationreach{John.4.51}{John}{true}\chronologyannotationreach{John.4.52}{John}{true}\chronologyannotationreach{John.4.53}{John}{true}\chronologyannotationclaim{composition.gospel-of-john}{composition}{catholic-traditional-v1}{disputed}{from the year 90 to 100 (approximately)}{c. A.D. 90--100}; \chronologyannotationreach{John.4.46}{John}{true}\chronologyannotationreach{John.4.47}{John}{true}\chronologyannotationreach{John.4.48}{John}{true}\chronologyannotationreach{John.4.49}{John}{true}\chronologyannotationreach{John.4.50}{John}{true}\chronologyannotationreach{John.4.51}{John}{true}\chronologyannotationreach{John.4.52}{John}{true}\chronologyannotationreach{John.4.53}{John}{true}\chronologyannotationclaim{composition.gospel-of-john}{composition}{catholic-traditional-v1}{disputed}{96 or one of the succeeding years}{A.D. 96--100}.}%
}
% offertory: Offertory
\expandafter\def\csname triptychchronologyannotation@offertory\endcsname{%
  \chronologyannotationgroup{historical-setting}{preferred}{\textbf{Historical setting}: \chronologyannotationreach{Ps.136.1}{Ps.136}{true}\chronologyannotationclaim{israel.exile.psalm-136-captive-lament}{historical-setting}{catholic-traditional-v1}{preferred}{For the time of Jeremias, and the captivity of Babylon}{The Babylonian captives' lament, after Jerusalem's destruction under Sedecias (derived)}.}%
  \space \chronologyannotationgroup{historical-setting}{preferred}{\textbf{Historical setting}: \chronologyannotationreach{Ps.136.1}{Ps.136}{true}\chronologyannotationclaim{israel.exile.seventy-years}{historical-setting}{catholic-traditional-v1}{preferred}{seventy years}{The seventy years of servitude to Babylon, ``seventy years'' (duration)}.}%
  \space \chronologyannotationgroup{composition}{preferred}{\textbf{Composition}: \chronologyannotationreach{Ps.136.1}{Ps}{true}\chronologyannotationclaim{critical.psalms.latest-composition-boundary}{composition}{catholic-critical-v1}{preferred}{before the Maccabean period, around 165 B.C.; no individual psalm can be dated securely}{Before c. 165 B.C.}}%
  \space \chronologyannotationanchorgroup{israel.exile.psalm-136-captive-lament}{historical-setting}{catholic-traditional-v1}{For the time of Jeremias, and the captivity of Babylon}{israel.exile.third-captivity}{after}{context}{\textbf{Historical background: The third captivity of Juda: the destruction of Jerusalem under Sedecias}: Preferred: \chronologyannotationanchorclaim{israel.exile.third-captivity}{0}{catholic-traditional-v1}{preferred}{passage.george-leo-haydock.douay-rheims-with-haydock-commentary.2014-loreto-feeney-memorial.psalm-70-captivities-usher-chronology}{reported-excluded}{A.M. 3416}{Haydock reporting Ussher; A.M. is not converted}{A.M. 3416 (Haydock reporting Ussher; A.M. is not converted)}; alternatives: \chronologyannotationanchorclaim{israel.exile.third-captivity}{1}{catholic-traditional-v1}{alternate}{artifact.catholic-encyclopedia.volume-3.new-york-1908.newadvent-03315a-9cb437e2,artifact.catholic-encyclopedia.volume-14.new-york-1912.newadvent-14499a-82b12b6e}{traditional-catholic}{586 B.C.}{Reid and Meistermann, Catholic Encyclopedia}{B.C. 586 (Reid and Meistermann, Catholic Encyclopedia)}, \chronologyannotationanchorclaim{israel.exile.third-captivity}{2}{catholic-traditional-v1}{alternate}{artifact.catholic-encyclopedia.volume-8.new-york-1910.newadvent-08652a-189e1a8f}{traditional-catholic}{the fall of Jerusalem (587 B.C.)}{Schets, Catholic Encyclopedia}{B.C. 587 (Schets, Catholic Encyclopedia)}, \chronologyannotationanchorclaim{israel.exile.third-captivity}{3}{catholic-traditional-v1}{alternate}{artifact.catholic-encyclopedia.volume-8.new-york-1910.newadvent-08654a-645bba6c}{reported-traditional}{B.C. 588}{Petavius table reported by Sloet; retained transcription checked, printed-page verification still pending}{B.C. 588 (Petavius table reported by Sloet; retained transcription checked, printed-page verification still pending)}. The represented event is after this historical event. These dates belong to that historical event, not to the passage or its composition. Both events could occur in the same year.}%
}
% communion: Communion
\expandafter\def\csname triptychchronologyannotation@communion\endcsname{%
  \chronologyannotationgroup{traditional-attribution}{disputed}{\textbf{Traditional attribution} -- disputed: \chronologyannotationreach{Ps.118.49}{Ps.118}{true}\chronologyannotationreach{Ps.118.50}{Ps.118}{true}\chronologyannotationclaim{israel.monarchy.david-traditional-era}{traditional-attribution}{catholic-traditional-v1}{disputed}{reigned from 1055 to 1015 B.C.}{David (reign in the usual chronology), B.C. 1055--1015}.}%
  \space \chronologyannotationgroup{historical-setting}{research-pending}{\textbf{Historical setting} -- Occasion date unresolved.}%
  \space \chronologyannotationgroup{composition}{preferred}{\textbf{Composition}: \chronologyannotationreach{Ps.118.49}{Ps}{true}\chronologyannotationreach{Ps.118.50}{Ps}{true}\chronologyannotationclaim{critical.psalms.latest-composition-boundary}{composition}{catholic-critical-v1}{preferred}{before the Maccabean period, around 165 B.C.; no individual psalm can be dated securely}{Before c. 165 B.C.}}%
}
% presentation: orientation-v3
\newcommand{\chronologyorientationrow}[2]{%
  \ifcsname triptychchronologyorientation@#1\endcsname
    \csname triptychchronologyorientation@#1\endcsname{#2}%
  \else
    \PackageError{triptych}{No chronology orientation row for #1}{}%
  \fi}
\expandafter\def\csname triptychchronologyorientation@row-6a30c5b77b668f6f\endcsname#1{%
Alleluia: Ps. 107:2 (Vg.; Heb. 108) & Traditional attribution: David's reign; Hebron, then Jerusalem & Disputed: \chronologyannotationreach{Ps.107.2}{Ps.107 Ps.144}{true}\chronologyannotationclaim{israel.monarchy.david-traditional-era}{traditional-attribution}{catholic-traditional-v1}{disputed}{reigned from 1055 to 1015 B.C.}{B.C.~1055--1015}\\*[.1em]%
}
\expandafter\def\csname triptychchronologyorientation@row-d09d8f0c8cf478ed\endcsname#1{%
 & Historical setting & Occasion date unresolved\\*[.1em]%
}
\expandafter\def\csname triptychchronologyorientation@row-d09dc302655154be\endcsname#1{%
 & Composition: Psalms & \chronologyannotationreach{Ps.107.2}{Ps}{true}\chronologyannotationclaim{critical.psalms.latest-composition-boundary}{composition}{catholic-critical-v1}{preferred}{before the Maccabean period, around 165 B.C.; no individual psalm can be dated securely}{Before c.~165 B.C.}\\[.3em]%
}
\expandafter\def\csname triptychchronologyorientation@row-1d7ed2f3a0e974d5\endcsname#1{%
Introit: Ps. 118:1 (Vg.; Heb. 119) & Traditional attribution: David's reign; Hebron, then Jerusalem & Disputed: \chronologyannotationreach{Ps.118.1}{Ps.118}{true}\chronologyannotationclaim{israel.monarchy.david-traditional-era}{traditional-attribution}{catholic-traditional-v1}{disputed}{reigned from 1055 to 1015 B.C.}{B.C.~1055--1015}\\*[.1em]%
}
\expandafter\def\csname triptychchronologyorientation@row-4333b00acdc5ef5d\endcsname#1{%
 & Historical setting & Occasion date unresolved\\*[.1em]%
}
\expandafter\def\csname triptychchronologyorientation@row-44d4e3a6bd7e728f\endcsname#1{%
 & Composition: Psalms & \chronologyannotationreach{Ps.118.1}{Ps}{true}\chronologyannotationclaim{critical.psalms.latest-composition-boundary}{composition}{catholic-critical-v1}{preferred}{before the Maccabean period, around 165 B.C.; no individual psalm can be dated securely}{Before c.~165 B.C.}\\[.3em]%
}
\expandafter\def\csname triptychchronologyorientation@row-dedd1a8ba0827634\endcsname#1{%
Communion: Ps. 118:49--50 (Vg.; Heb. 119) & Traditional attribution: David's reign; Hebron, then Jerusalem & Disputed: \chronologyannotationreach{Ps.118.49}{Ps.118}{true}\chronologyannotationreach{Ps.118.50}{Ps.118}{true}\chronologyannotationclaim{israel.monarchy.david-traditional-era}{traditional-attribution}{catholic-traditional-v1}{disputed}{reigned from 1055 to 1015 B.C.}{B.C.~1055--1015}\\*[.1em]%
}
\expandafter\def\csname triptychchronologyorientation@row-a397348542bdd78f\endcsname#1{%
 & Historical setting & Occasion date unresolved\\*[.1em]%
}
\expandafter\def\csname triptychchronologyorientation@row-21e0d1248576c2af\endcsname#1{%
 & Composition: Psalms & \chronologyannotationreach{Ps.118.49}{Ps}{true}\chronologyannotationreach{Ps.118.50}{Ps}{true}\chronologyannotationclaim{critical.psalms.latest-composition-boundary}{composition}{catholic-critical-v1}{preferred}{before the Maccabean period, around 165 B.C.; no individual psalm can be dated securely}{Before c.~165 B.C.}\\[.3em]%
}
\expandafter\def\csname triptychchronologyorientation@row-ce9f6691f9a565f5\endcsname#1{%
Offertory: Ps. 136:1 (Vg.; Heb. 137) & Historical setting: The Babylonian captives' lament & \chronologyannotationreach{Ps.136.1}{Ps.136}{true}\chronologyannotationclaim{israel.exile.psalm-136-captive-lament}{historical-setting}{catholic-traditional-v1}{preferred}{For the time of Jeremias, and the captivity of Babylon}{After Jerusalem's destruction under Sedecias (derived)}\\*[.1em]%
}
\expandafter\def\csname triptychchronologyorientation@row-52147c6d72df9b71\endcsname#1{%
 & Historical setting: The seventy years of servitude to Babylon; Judah and Babylon & \chronologyannotationreach{Ps.136.1}{Ps.136}{true}\chronologyannotationclaim{israel.exile.seventy-years}{historical-setting}{catholic-traditional-v1}{preferred}{seventy years}{``seventy years'' (duration)}\\*[.1em]%
}
\expandafter\def\csname triptychchronologyorientation@row-84005e42520c6919\endcsname#1{%
 & Composition: Psalms & \chronologyannotationreach{Ps.136.1}{Ps}{true}\chronologyannotationclaim{critical.psalms.latest-composition-boundary}{composition}{catholic-critical-v1}{preferred}{before the Maccabean period, around 165 B.C.; no individual psalm can be dated securely}{Before c.~165 B.C.}\\*[.1em]%
}
\expandafter\def\csname triptychchronologyorientation@row-02760c799c7897ae\endcsname#1{%
 & Historical background: Jerusalem's destruction under Sedecias; setting after this event, possibly in the same year & Preferred: \chronologyannotationanchorclaim{israel.exile.third-captivity}{0}{catholic-traditional-v1}{preferred}{passage.george-leo-haydock.douay-rheims-with-haydock-commentary.2014-loreto-feeney-memorial.psalm-70-captivities-usher-chronology}{reported-excluded}{A.M. 3416}{Haydock reporting Ussher; A.M. is not converted}{A.M.~3416 (Haydock reporting Ussher)}; alternatives: \chronologyannotationanchorclaim{israel.exile.third-captivity}{1}{catholic-traditional-v1}{alternate}{artifact.catholic-encyclopedia.volume-3.new-york-1908.newadvent-03315a-9cb437e2,artifact.catholic-encyclopedia.volume-14.new-york-1912.newadvent-14499a-82b12b6e}{traditional-catholic}{586 B.C.}{Reid and Meistermann, Catholic Encyclopedia}{B.C.~586 (Reid and Meistermann)}, \chronologyannotationanchorclaim{israel.exile.third-captivity}{2}{catholic-traditional-v1}{alternate}{artifact.catholic-encyclopedia.volume-8.new-york-1910.newadvent-08652a-189e1a8f}{traditional-catholic}{the fall of Jerusalem (587 B.C.)}{Schets, Catholic Encyclopedia}{B.C.~587 (Schets)}, \chronologyannotationanchorclaim{israel.exile.third-captivity}{3}{catholic-traditional-v1}{alternate}{artifact.catholic-encyclopedia.volume-8.new-york-1910.newadvent-08654a-645bba6c}{reported-traditional}{B.C. 588}{Petavius table reported by Sloet; retained transcription checked, printed-page verification still pending}{B.C.~588 (Petavius via Sloet; transcription only)}\\[.3em]%
}
\expandafter\def\csname triptychchronologyorientation@row-51f8eb0674af573d\endcsname#1{%
Gradual: Ps. 144:15--16 (Vg.; Heb. 145) & Traditional attribution: David's reign; Hebron, then Jerusalem & Disputed: \chronologyannotationreach{Ps.144.15}{Ps.107 Ps.144}{true}\chronologyannotationreach{Ps.144.16}{Ps.107 Ps.144}{true}\chronologyannotationclaim{israel.monarchy.david-traditional-era}{traditional-attribution}{catholic-traditional-v1}{disputed}{reigned from 1055 to 1015 B.C.}{B.C.~1055--1015}\\*[.1em]%
}
\expandafter\def\csname triptychchronologyorientation@row-56fc764f3317104e\endcsname#1{%
 & Historical setting & Occasion date unresolved\\*[.1em]%
}
\expandafter\def\csname triptychchronologyorientation@row-691db73eda44086c\endcsname#1{%
 & Composition: Psalms & \chronologyannotationreach{Ps.144.15}{Ps}{true}\chronologyannotationreach{Ps.144.16}{Ps}{true}\chronologyannotationclaim{critical.psalms.latest-composition-boundary}{composition}{catholic-critical-v1}{preferred}{before the Maccabean period, around 165 B.C.; no individual psalm can be dated securely}{Before c.~165 B.C.}\\[.3em]%
}
\expandafter\def\csname triptychchronologyorientation@row-a88ec9dc1b6cbc4c\endcsname#1{%
Introit: Dan. 3:29,31,35 & Historical setting: The seventy years of servitude to Babylon; Judah and Babylon & \chronologyannotationreach{Dan.3.29}{Dan}{true}\chronologyannotationreach{Dan.3.31}{Dan}{true}\chronologyannotationreach{Dan.3.35}{Dan}{true}\chronologyannotationclaim{israel.exile.seventy-years}{historical-setting}{catholic-traditional-v1}{preferred}{seventy years}{``seventy years'' (duration)}\\*[.1em]%
}
\expandafter\def\csname triptychchronologyorientation@row-29eede98fe3a316e\endcsname#1{%
 & Composition: Prayer of Azarias and Song of the Three Children & \chronologyannotationreach{Dan.3.29}{Dan.3.24-90}{false}\chronologyannotationreach{Dan.3.31}{Dan.3.24-90}{false}\chronologyannotationreach{Dan.3.35}{Dan.3.24-90}{false}\chronologyannotationclaim{composition.book-of-daniel.prayer-of-azarias}{composition}{catholic-traditional-v1}{preferred}{most likely long after the Exile}{Most likely long after the Exile}\\[.3em]%
}
\expandafter\def\csname triptychchronologyorientation@row-5ebc75edfb8977e6\endcsname#1{%
Gospel: John 4:46--53 & Narrated event: The healing of the ruler's son at Cana & \chronologyannotationreach{John.4.46}{John.4.46-53}{false}\chronologyannotationreach{John.4.47}{John.4.46-53}{false}\chronologyannotationreach{John.4.48}{John.4.46-53}{false}\chronologyannotationreach{John.4.49}{John.4.46-53}{false}\chronologyannotationreach{John.4.50}{John.4.46-53}{false}\chronologyannotationreach{John.4.51}{John.4.46-53}{false}\chronologyannotationreach{John.4.52}{John.4.46-53}{false}\chronologyannotationreach{John.4.53}{John.4.46-53}{false}\chronologyannotationclaim{life-of-christ.healing-of-the-rulers-son-at-cana}{narrated-event}{catholic-traditional-v1}{preferred}{Passover, A.U.C. 779 - about Pentecost, 780}{A.D.~26--27 (derived)}\\*[.1em]%
}
\expandafter\def\csname triptychchronologyorientation@row-f3e7eb0ac2152a88\endcsname#1{%
 & Composition: John; Ephesus traditionally & Disputed: \chronologyannotationreach{John.4.46}{John}{true}\chronologyannotationreach{John.4.47}{John}{true}\chronologyannotationreach{John.4.48}{John}{true}\chronologyannotationreach{John.4.49}{John}{true}\chronologyannotationreach{John.4.50}{John}{true}\chronologyannotationreach{John.4.51}{John}{true}\chronologyannotationreach{John.4.52}{John}{true}\chronologyannotationreach{John.4.53}{John}{true}\chronologyannotationclaim{composition.gospel-of-john}{composition}{catholic-traditional-v1}{disputed}{from the year 90 to 100 (approximately)}{c.~A.D.~90--100}; \chronologyannotationreach{John.4.46}{John}{true}\chronologyannotationreach{John.4.47}{John}{true}\chronologyannotationreach{John.4.48}{John}{true}\chronologyannotationreach{John.4.49}{John}{true}\chronologyannotationreach{John.4.50}{John}{true}\chronologyannotationreach{John.4.51}{John}{true}\chronologyannotationreach{John.4.52}{John}{true}\chronologyannotationreach{John.4.53}{John}{true}\chronologyannotationclaim{composition.gospel-of-john}{composition}{catholic-traditional-v1}{disputed}{96 or one of the succeeding years}{A.D.~96--100}\\[.3em]%
}
\expandafter\def\csname triptychchronologyorientation@row-3db431fd3c4a86c3\endcsname#1{%
Epistle: Eph. 5:15--21 & Composition: Ephesians; Rome or Caesarea; recipients in Asia Minor & Disputed: \chronologyannotationreach{Eph.5.15}{Eph}{true}\chronologyannotationreach{Eph.5.16}{Eph}{true}\chronologyannotationreach{Eph.5.17}{Eph}{true}\chronologyannotationreach{Eph.5.18}{Eph}{true}\chronologyannotationreach{Eph.5.19}{Eph}{true}\chronologyannotationreach{Eph.5.20}{Eph}{true}\chronologyannotationreach{Eph.5.21}{Eph}{true}\chronologyannotationclaim{composition.epistle-to-the-ephesians}{composition}{catholic-traditional-v1}{disputed}{a period between 58 and 63}{A.D.~58--63}; \chronologyannotationreach{Eph.5.15}{Eph}{true}\chronologyannotationreach{Eph.5.16}{Eph}{true}\chronologyannotationreach{Eph.5.17}{Eph}{true}\chronologyannotationreach{Eph.5.18}{Eph}{true}\chronologyannotationreach{Eph.5.19}{Eph}{true}\chronologyannotationreach{Eph.5.20}{Eph}{true}\chronologyannotationreach{Eph.5.21}{Eph}{true}\chronologyannotationclaim{composition.epistle-to-the-ephesians}{composition}{catholic-traditional-v1}{disputed}{(Philemon; Colossians; Ephesians; Philippians), 61}{A.D.~61}\\*[.1em]%
}
\expandafter\def\csname triptychchronologyorientation@row-d253ee44aa17d758\endcsname#1{%
 & Composition, Catholic critical comparison: Ephesians (later-disciple hypothesis) & Disputed: \chronologyannotationreach{Eph.5.15}{Eph}{true}\chronologyannotationreach{Eph.5.16}{Eph}{true}\chronologyannotationreach{Eph.5.17}{Eph}{true}\chronologyannotationreach{Eph.5.18}{Eph}{true}\chronologyannotationreach{Eph.5.19}{Eph}{true}\chronologyannotationreach{Eph.5.20}{Eph}{true}\chronologyannotationreach{Eph.5.21}{Eph}{true}\chronologyannotationcomparisonclaim{critical.ephesians-later-disciple}{composition}{catholic-critical-v1}{catholic-critical-v1}{disputed}{passage.united-states-conference-of-catholic-bishops.new-american-bible-revised-edition.english-usccb-web-2026-09-28.ephesians-introduction}{around A.D. 80--100}{c.~A.D.~80--100}\\[.3em]%
}

% Generated. Do not edit.
% Complete orientation-v3 table; wording, dates and order belong to the corpus.
\begin{chronologysharedtable}
\chronologyorientationrow{row-6a30c5b77b668f6f}{}
\chronologyorientationrow{row-d09d8f0c8cf478ed}{}
\chronologyorientationrow{row-d09dc302655154be}{}
\midrule
\chronologyorientationrow{row-1d7ed2f3a0e974d5}{}
\chronologyorientationrow{row-4333b00acdc5ef5d}{}
\chronologyorientationrow{row-44d4e3a6bd7e728f}{}
\midrule
\chronologyorientationrow{row-dedd1a8ba0827634}{}
\chronologyorientationrow{row-a397348542bdd78f}{}
\chronologyorientationrow{row-21e0d1248576c2af}{}
\midrule
\chronologyorientationrow{row-ce9f6691f9a565f5}{}
\chronologyorientationrow{row-52147c6d72df9b71}{}
\chronologyorientationrow{row-84005e42520c6919}{}
\chronologyorientationrow{row-02760c799c7897ae}{}
\midrule
\chronologyorientationrow{row-51f8eb0674af573d}{}
\chronologyorientationrow{row-56fc764f3317104e}{}
\chronologyorientationrow{row-691db73eda44086c}{}
\midrule
\chronologyorientationrow{row-a88ec9dc1b6cbc4c}{}
\chronologyorientationrow{row-29eede98fe3a316e}{}
\midrule
\chronologyorientationrow{row-5ebc75edfb8977e6}{}
\chronologyorientationrow{row-f3e7eb0ac2152a88}{}
\midrule
\chronologyorientationrow{row-3db431fd3c4a86c3}{}
\chronologyorientationrow{row-d253ee44aa17d758}{}
\end{chronologysharedtable}
{\footnotesize Psalm numbering: Vg. = Vulgate; Heb. = Hebrew.\par}


\subsection{Evidence for each date}
\label{sec:date-location-evidence}
{\small For each passage, as the Scriptural Date and Location table names it: what is dated, how the corpus weighs each answer, the source's own words, and the sources cited. A subject the corpus does not date is reported as such and given no date.\par}
{\small
\par\smallskip\noindent\textit{Alleluia: Ps. 107:2 (Vg.; Heb. 108)}\par\nobreak
\begin{itemize}[leftmargin=1.2em,itemsep=1pt,topsep=1pt,parsep=0pt]
\item \textbf{Traditional attribution} (disputed). David's reign: B.C.~1055--1015 (``reigned from 1055 to 1015 B.C.''; The Catholic Encyclopedia, Volume 4 (New York, 1908), article 'David, King,' printed p. 642, opening paragraph). Place: Hebron, then Jerusalem (The Holy Bible, Douay-Rheims (Challoner revision), Project Gutenberg eBook 1581: 2 Kings 5:5).
\item \textbf{Historical setting}. No date: the chronology corpus has not yet dated this (research pending).
\item \textbf{Composition} (preferred). Psalms: Before c.~165 B.C. (``before the Maccabean period, around 165 B.C.; no individual psalm can be dated securely''; Catholic critical chronology; New American Bible, Revised Edition --- USCCB web state (2026), Psalms introduction).
\end{itemize}
\par\smallskip\noindent\textit{Introit: Ps. 118:1 (Vg.; Heb. 119)}\par\nobreak
\begin{itemize}[leftmargin=1.2em,itemsep=1pt,topsep=1pt,parsep=0pt]
\item \textbf{Traditional attribution}. As under Alleluia: Ps. 107:2 (Vg.; Heb. 108).
\item \textbf{Historical setting}. As under Alleluia: Ps. 107:2 (Vg.; Heb. 108).
\item \textbf{Composition}. As under Alleluia: Ps. 107:2 (Vg.; Heb. 108).
\end{itemize}
\par\smallskip\noindent\textit{Communion: Ps. 118:49--50 (Vg.; Heb. 119)}\par\nobreak
\begin{itemize}[leftmargin=1.2em,itemsep=1pt,topsep=1pt,parsep=0pt]
\item \textbf{Traditional attribution}. As under Alleluia: Ps. 107:2 (Vg.; Heb. 108).
\item \textbf{Historical setting}. As under Alleluia: Ps. 107:2 (Vg.; Heb. 108).
\item \textbf{Composition}. As under Alleluia: Ps. 107:2 (Vg.; Heb. 108).
\end{itemize}
\par\smallskip\noindent\textit{Offertory: Ps. 136:1 (Vg.; Heb. 137)}\par\nobreak
\begin{itemize}[leftmargin=1.2em,itemsep=1pt,topsep=1pt,parsep=0pt]
\item \textbf{Historical setting} (preferred). The Babylonian captives' lament: After Jerusalem's destruction under Sedecias (derived) (``For the time of Jeremias, and the captivity of Babylon''; The Holy Bible, Douay-Rheims (Challoner revision), Project Gutenberg eBook 1581: Ps. 136:1, 136:3, 136:7; Douay Rheims Edition of the Holy Catholic Bible with a Comprehensive Catholic Commentary, William Peter Feeney 2014 Memorial Edition, Ps.136.1).
\item \textbf{Historical setting} (preferred). The seventy years of servitude to Babylon: ``seventy years'' (duration) (The Holy Bible, Douay-Rheims (Challoner revision), Project Gutenberg eBook 1581: Jer. 25:11, 25:12). Place: Judah and Babylon (The Holy Bible, Douay-Rheims (Challoner revision), Project Gutenberg eBook 1581: Jer. 25:11, 25:12).
\item \textbf{Composition}. As under Alleluia: Ps. 107:2 (Vg.; Heb. 108).
\item \textbf{Historical background} (answers differ). Jerusalem's destruction under Sedecias: preferred: A.M.~3416 (Haydock reporting Ussher) (Douay Rheims Edition of the Holy Catholic Bible with a Comprehensive Catholic Commentary, William Peter Feeney 2014 Memorial Edition, Ps.70.1); alternative: B.C.~586 (Reid and Meistermann) (``586 B.C.''; The Catholic Encyclopedia, Volume 3 (New York, 1908); The Catholic Encyclopedia, Volume 14 (New York, 1912)); alternative: B.C.~587 (Schets) (``the fall of Jerusalem (587 B.C.)''; The Catholic Encyclopedia, Volume 8 (New York, 1910)); alternative: B.C.~588 (Petavius via Sloet; transcription only) (The Catholic Encyclopedia, Volume 8 (New York, 1910)). These dates belong to that event, not to the passage or its composition, and the two could fall in the same year.
\end{itemize}
\par\smallskip\noindent\textit{Gradual: Ps. 144:15--16 (Vg.; Heb. 145)}\par\nobreak
\begin{itemize}[leftmargin=1.2em,itemsep=1pt,topsep=1pt,parsep=0pt]
\item \textbf{Traditional attribution}. As under Alleluia: Ps. 107:2 (Vg.; Heb. 108).
\item \textbf{Historical setting}. As under Alleluia: Ps. 107:2 (Vg.; Heb. 108).
\item \textbf{Composition}. As under Alleluia: Ps. 107:2 (Vg.; Heb. 108).
\end{itemize}
\par\smallskip\noindent\textit{Introit: Dan. 3:29,31,35}\par\nobreak
\begin{itemize}[leftmargin=1.2em,itemsep=1pt,topsep=1pt,parsep=0pt]
\item \textbf{Historical setting}. As under Offertory: Ps. 136:1 (Vg.; Heb. 137).
\item \textbf{Composition} (preferred). Prayer of Azarias and Song of the Three Children: Most likely long after the Exile (The Catholic Encyclopedia, Volume 4 (New York, 1908)).
\end{itemize}
\par\smallskip\noindent\textit{Gospel: John 4:46--53}\par\nobreak
\begin{itemize}[leftmargin=1.2em,itemsep=1pt,topsep=1pt,parsep=0pt]
\item \textbf{Narrated event} (preferred). The healing of the ruler's son at Cana: A.D.~26--27 (derived) (``Passover, A.U.C. 779 - about Pentecost, 780''; The Catholic Encyclopedia, Volume 8 (New York, 1910), Second journey: Cana and the dying son at Capharnaum; source-internal chronological calibration).
\item \textbf{Composition} (disputed). The Gospel of St John: c.~A.D.~90--100 (``from the year 90 to 100 (approximately)''; The Catholic Encyclopedia, Volume 14 (New York, 1912)); A.D.~96--100 (``96 or one of the succeeding years''; The Catholic Encyclopedia, Volume 8 (New York, 1910)). Place: Ephesus traditionally (The Catholic Encyclopedia, Volume 8 (New York, 1910)).
\end{itemize}
\par\smallskip\noindent\textit{Epistle: Eph. 5:15--21}\par\nobreak
\begin{itemize}[leftmargin=1.2em,itemsep=1pt,topsep=1pt,parsep=0pt]
\item \textbf{Composition} (disputed). The Epistle to the Ephesians: A.D.~58--63 (``a period between 58 and 63''; The Catholic Encyclopedia, Volume 5 (New York, 1909)); A.D.~61 (``(Philemon; Colossians; Ephesians; Philippians), 61''; The Catholic Encyclopedia, Volume 11 (New York, 1911)). Place: Rome or Caesarea; recipients in Asia Minor (The Catholic Encyclopedia, Volume 5 (New York, 1909)).
\item \textbf{Composition, Catholic critical comparison} (disputed). Ephesians under the later-disciple hypothesis (approximate): c.~A.D.~80--100 (``around A.D. 80--100''; Catholic critical chronology; New American Bible, Revised Edition: USCCB web state (2026), The Letter to the Ephesians, complete introduction; especially prose paragraphs 3-4, physical HTML lines 626 and 628).
\end{itemize}
\par}
